\usepackage[utf8]{inputenc} %Eingabekodierung
\usepackage[T1]{fontenc} %Silbentrennung bei Sonderzeichen

\usepackage{graphicx} %Bilder einbinden
\usepackage[ngerman]{babel} %Deutsche Sprachanpassungen
\usepackage{csquotes} %Deutsche Anführungszeichen
\usepackage{acronym} %für optionales Abkürzungsverzeichnis
\usepackage{eurosym} %z. B. \EUR{12345,68}
%\usepackage[bindingoffset=8mm]{geometry} %Bindeverlust von 8mm einbeziehen. Mit dem geometry-Paket können Sie die Ränder auch ganz individuell anpassen.
\usepackage[left=1in,right=1in,top=1in,bottom=1in]{geometry}
%\usepackage[a4paper,left=25mm,right=25mm,top=25mm,bottom=25mm]{geometry} %Zum zentrieren der Seite (muss aber dann bindingoffset deaktiviert werden)
%TODO:prüfen ob benötigt \usepackage[a-2b,mathxmp]{pdfx} %PDF/A-2b Konformität
\usepackage{caption} %Abbildungslegenden
\usepackage{float} %Zum Plazieren von z.B. Tabellen

%Schriften
% Source https://tex.stackexchange.com/a/507509
\usepackage[onehalfspacing]{setspace}
% set font to helvetica
\usepackage{helvet}
\renewcommand*{\familydefault}{\sfdefault}


%Optisch
\usepackage{pdflscape} %Querformat-Seiten
\usepackage{listings} %Codeanzeige
\usepackage{enumitem} %Ermöglicht das anpassen von Listen wie (itemize)
\usepackage{siunitx} %Zahlenausrichtung
\usepackage[hyphens]{xurl} %intelligente Umbrüche in URLs

% Customize chapter style
\RedeclareSectionCommand[%
  innerskip=0.05\baselineskip,
  afterskip=1\baselineskip plus .1\baselineskip minus .12\baselineskip,
  beforeskip=0pt,       % Removes the top padding completely
  afterindent=false     % Modern KOMA way to suppress paragraph indent on the next line
]{chapter}

%Customize section and subsection distances to match chapter style
\RedeclareSectionCommands[
runin=false,              %% NB! Important
afterskip=-0.5\parskip
]{section,subsection}

%Literaturverweise
\usepackage[hidelinks]{hyperref} %sorgt für klickbare links von Zitaten und von Links im literaturverzeichnis
\usepackage[style=apa,citestyle=apa,backend=biber,hyperref=true,doi=true,url=true]{biblatex} %Literaturverweise
%biblatex comes with a variety of built-in bibliography/citation style families (numeric, alphabetic, authoryear, authortitle, verbose), and there's a growing number of custom styles:
%https://de.sharelatex.com/learn/Biblatex_citation_styles
%https://de.sharelatex.com/learn/Biblatex_bibliography_styles

%Tabellen
\usepackage{booktabs} %Professionelle Tabellenlinien
\usepackage{longtable} %Tabellen über mehrere Seiten
\usepackage{makecell} %Spalten in Tabellen
\usepackage[table]{xcolor}  %Option 'table' für Tabellenfarben
\usepackage{tabularx} %Automatische Breitenanpassung
\usepackage{array} %Verbesserte Spaltenformatierung
%\usepackage{array} %Zentrierte text in spalte